\chapter{Conclusiones}\label{cap:conclusiones}

Esta investigación modeló la movilidad de Bogotá y su región con una demanda estocástica, un movimiento
browniano geométrico multivariado, y evaluó con simulación de Monte Carlo los diseños urbanos que
combinan inversiones férreas y política de suelo. El trabajo se desarrolló en dos fases: la primera con
la coyuntura de 2013 y la segunda con las condiciones de 2025. Las conclusiones principales son las
siguientes.

\paragraph{Metodología.} La solución exacta del MBG, los números aleatorios comunes y el cuasi-Monte
Carlo con secuencias de Sobol permiten ordenar decenas de diseños con errores de muestreo muy pequeños
frente a las diferencias entre ellos: el error cuadrático medio del costo esperado converge a una tasa
cercana a $n^{\PendienteSobol}$, frente a $n^{\PendienteMC}$ del Monte Carlo simple. El algoritmo de
mínimos cuadrados de Monte Carlo reproduce la solución analítica de McDonald y Siegel con un error de
\MSError.

\paragraph{Fase I: 2013.} Con la información de 2013 y la tasa de descuento de 12\,\%, el metro elevado
tenía el mayor beneficio neto esperado entre las alternativas de primera línea (\BNElevDoce{} billones,
frente a \BNSubOfDoce{} del subterráneo con su costo oficial). La demanda que efectivamente ocurrió
quedó en el percentil \PercentilObsFin{} de los escenarios de 2013, por la menor población, la
pandemia y una recuperación incompleta; con ella, ninguna primera línea se paga a 12\,\%, y el
subterráneo habría perdido más que el elevado.

\paragraph{Fase II: 2025.} Con la tasa de 9\,\%, el diseño óptimo es \OptimoNeutral: el Regiotram del
Norte con la política de proximidad completa, con un beneficio neto esperado de \BNOptimo{} billones y
positivo en el \PBNOptimo{} de los escenarios. Las líneas de metro adicionales reducen la dispersión
del costo social pero no se pagan a esa tasa; con saltos sistémicos, la combinación con la Línea 3
tiene el menor CVaR. La Línea 3 tiene un valor de opción pequeño pero positivo (\OpcionFlexible{}
billones), frente a un VPN de \OpcionAhora{} si se compromete hoy: la regla es esperar.

\paragraph{El metro de Bogotá.} La hipótesis de que la Primera Línea debió ser subterránea por el tamaño
de la ciudad no se sostiene en su forma general, pero tiene un núcleo defendible. El tamaño, por sí
solo, no la justifica: la capacidad del viaducto no es la restricción y una ciudad grande necesita una
red, para la cual el elevado permite construir \KmPorKm{} kilómetros por cada kilómetro subterráneo.
El subterráneo de 2014 habría sido preferible con una tasa de descuento inferior a \TasaEmpate{}, o con
la tasa actual de 9\,\% si el costo urbano del viaducto supera unos \UmbralOfNueve{} billones por año,
un valor que está en el extremo alto del rango estimado por pérdida de valor de los predios. Con la tasa
de la época, ese costo tendría que haber superado \UmbralOfDoce{} billones por año. Para las líneas 2 y
3, la tipología subterránea cuesta entre \AhorroLdos{} y \AhorroLtres{} billones más que una variante
elevada y exige una justificación explícita del corredor; una Línea 3 elevada entraría en el diseño
óptimo. Reabrir hoy la tipología de la Primera Línea no tiene sentido económico.

\paragraph{Recomendaciones.}
\begin{enumerate}[nosep]
\item Evaluar los proyectos de transporte sobre distribuciones de escenarios y reportar, además del
      valor esperado, el CVaR y la probabilidad de beneficio neto positivo.
\item Presentar las decisiones de tipología con su tasa de descuento de quiebre y el costo urbano
      umbral, de modo que el debate se haga sobre supuestos explícitos.
\item Priorizar la política de proximidad y los proyectos de menor costo por kilómetro antes que nuevas
      líneas subterráneas, y mantener la Línea 3 como opción, con una regla de ejercicio explícita.
\item Medir el costo urbano real del viaducto de la Primera Línea cuando entre en operación, para que
      las decisiones de las próximas líneas se apoyen en evidencia local.
\end{enumerate}

Todo el modelo, los datos y la aplicación interactiva están disponibles en el repositorio público, de
modo que cada cifra de este documento puede reproducirse y cada supuesto puede cambiarse.
